% =====================================================================
\chapter{Derivadas parciales y gradiente}
\label{tema:1}
% =====================================================================

La pérdida de un modelo de lenguaje es una función escalar de muchas variables. Este tema define las derivadas parciales y la diferenciabilidad, introduce el gradiente y la derivada direccional, y demuestra que el gradiente señala la dirección de máximo ascenso, que es el fundamento del descenso de gradiente. Se siguen las clases~7, 10 y~11 del curso 18.02 del MIT~\cite{mit1802}, la §5.2 de~\cite{mml} y la §2.4 de~\cite{d2l}.

\section{Derivadas parciales}

Sea $U\subseteq\R^{n}$ un abierto y $f\colon U\to\R$. Para estudiar cómo varía $f$ cerca de un punto, se empieza por moverse en una sola coordenada.

\begin{definicion}[Derivada parcial]\label{def:parcial}
La \emph{derivada parcial} de $f$ respecto de $x_i$ en el punto $x\in U$ es
\[
\frac{\partial f}{\partial x_i}(x)=\lim_{h\to 0}\frac{f(x+h\,e_i)-f(x)}{h},
\]
si el límite existe. Equivale a derivar $f$ como función de una sola variable $x_i$, tratando las demás como constantes.
\end{definicion}

\addtocounter{definicion}{1}%

\addtocounter{definicion}{1}%

\section{Diferenciabilidad y gradiente}

La existencia de las derivadas parciales no basta para que $f$ se comporte bien cerca de $x$: las parciales solo exploran las direcciones de los ejes. La noción adecuada exige que $f$ admita una buena aproximación lineal en todas las direcciones a la vez~\cite[§5.2]{mml}.

\begin{definicion}[Diferenciabilidad y gradiente]\label{def:diferenciable}
La función $f$ es \emph{diferenciable} en $x\in U$ si existe un vector $g\in\R^{n}$ tal que
\[
f(x+h)=f(x)+g\T h+o(\norm{h})\qquad\text{cuando } h\to 0,
\]
es decir, $\lim_{h\to 0}\bigl(f(x+h)-f(x)-g\T h\bigr)/\norm{h}=0$. El vector $g$ se llama \emph{gradiente} de $f$ en $x$ y se escribe $\nabla f(x)$. La función afín $h\mapsto f(x)+\nabla f(x)\T h$ es la \emph{aproximación lineal} (o de primer orden) de $f$ en $x$.
\end{definicion}

\begin{proposicion}\label{prop:grad-parciales}
Si $f$ es diferenciable en $x$, existen todas las parciales en $x$ y
\[
\nabla f(x)=\Bigl(\frac{\partial f}{\partial x_1}(x),\dots,\frac{\partial f}{\partial x_n}(x)\Bigr)\T.
\]
En particular, el gradiente es único.
\end{proposicion}

El gradiente tiene la misma forma que la variable: si $x\in\R^{n}$, entonces $\nabla f(x)\in\R^{n}$. En el \cref{ej:parciales}, $\nabla f(2,1)=(4,7)\T$. Para un LLM, el gradiente de la pérdida tiene tantas componentes como parámetros.

\begin{teorema}[Condición suficiente de diferenciabilidad]\label{teo:C1}
Si todas las derivadas parciales de $f$ existen en un entorno de $x$ y son continuas en $x$, entonces $f$ es diferenciable en $x$. En particular, toda $f\in C^{1}(U)$ es diferenciable en todo punto de $U$.
\end{teorema}

La demostración, estándar en los cursos de cálculo de varias variables, aplica el teorema del valor medio en cada coordenada por separado y se omite aquí. Las funciones que se usan para construir una red (polinomios, exponenciales, logaritmos, cocientes con denominador no nulo y sus composiciones) son de clase $C^{\infty}$ en su dominio. La excepción habitual es la función ReLU, $\max(0,t)$, que no es derivable en $t=0$; el \cref{tema:6} vuelve sobre ella.

\begin{observacion}[Parciales sin diferenciabilidad]
Sea $f(x,y)=xy/(x^{2}+y^{2})$ si $(x,y)\ne(0,0)$ y $f(0,0)=0$. Como $f$ se anula sobre los dos ejes, $\partial f/\partial x(0,0)=\partial f/\partial y(0,0)=0$. Sin embargo, $f(t,t)=1/2$ para todo $t\ne 0$, de modo que $f$ ni siquiera es continua en el origen y, por tanto, no es diferenciable allí. Las parciales existen, pero no son continuas en $(0,0)$.
\end{observacion}

\section{Derivada direccional}

\begin{definicion}[Derivada direccional]\label{def:direccional}
Sea $u\in\R^{n}$ con $\norm{u}=1$. La \emph{derivada direccional} de $f$ en $x$ en la dirección $u$ es
\[
D_u f(x)=\lim_{t\to 0}\frac{f(x+t\,u)-f(x)}{t}.
\]
Las derivadas parciales son las derivadas direccionales en las direcciones $e_1,\dots,e_n$.
\end{definicion}

\begin{proposicion}\label{prop:direccional}
Si $f$ es diferenciable en $x$, entonces $D_u f(x)=\nabla f(x)\T u$ para todo $u$ unitario.
\end{proposicion}

\addtocounter{definicion}{1}%

\section{La dirección de máximo ascenso}

\begin{teorema}[Máximo ascenso]\label{teo:max-ascenso}
Sea $f$ diferenciable en $x$ con $\nabla f(x)\ne 0$. Para todo $u$ unitario,
\[
-\norm{\nabla f(x)}\le D_u f(x)\le\norm{\nabla f(x)}.
\]
La cota superior se alcanza únicamente en $u^{*}=\nabla f(x)/\norm{\nabla f(x)}$ y la inferior únicamente en $-u^{*}$.
\end{teorema}

Así, $\nabla f(x)$ apunta en la dirección en la que $f$ crece más deprisa, su norma es esa tasa máxima de crecimiento y $-\nabla f(x)$ es la dirección de \emph{máximo descenso}.

\begin{lema}[Derivada a lo largo de una curva]\label{lem:curva}
Sea $r\colon(a,b)\to U$ derivable en $t_0$ y $f$ diferenciable en $r(t_0)$. Entonces $t\mapsto f(r(t))$ es derivable en $t_0$ y
\[
\frac{d}{dt}f(r(t))\Big|_{t=t_0}=\nabla f(r(t_0))\T r'(t_0).
\]
\end{lema}

El \cref{lem:curva} es un caso particular de la regla de la cadena, que se estudia en general en el \cref{tema:2}.

\begin{corolario}[Gradiente y conjuntos de nivel]\label{cor:nivel}
Sea $r$ una curva derivable contenida en un conjunto de nivel de $f$, es decir, $f(r(t))=c$ para todo $t$. Entonces $\nabla f(r(t))\T r'(t)=0$: el gradiente es ortogonal a la recta tangente a la curva de nivel.
\end{corolario}

\addtocounter{figure}{1}%

\section{Descenso de gradiente}

El \cref{teo:max-ascenso} sugiere un método para minimizar $f$: desplazarse repetidamente en la dirección $-\nabla f$.

\begin{definicion}[Descenso de gradiente]\label{def:dg}
Dado un punto inicial $x_0$ y una \emph{tasa de aprendizaje} $\eta>0$, el \emph{descenso de gradiente} genera la sucesión
\[
x_{k+1}=x_k-\eta\,\nabla f(x_k),\qquad k=0,1,2,\dots
\]
\end{definicion}

\begin{proposicion}[Descenso a primer orden]\label{prop:descenso}
Si $f$ es diferenciable en $x$ y $g=\nabla f(x)\ne 0$, existe $\eta_0>0$ tal que $f(x-\eta g)<f(x)$ para todo $\eta\in(0,\eta_0)$.
\end{proposicion}

La proposición garantiza el descenso solo para pasos pequeños. El tamaño máximo admisible depende de la curvatura de $f$ y se determina en el \cref{tema:4}.

\addtocounter{definicion}{1}%

\addtocounter{definicion}{1}%

\begin{observacion}[Relación con la física]
Una partícula en un potencial $U(x)$ sometida a un rozamiento muy fuerte se mueve con velocidad proporcional a la fuerza $-\nabla U$: $\dot x(t)=-\nabla U(x(t))$. Esta ecuación se llama \emph{flujo de gradiente}, y la \cref{def:dg} es su discretización por el método de Euler con paso $\eta$. Añadir un término de ruido da la dinámica de Langevin, que sirve de modelo para el descenso de gradiente estocástico.
\end{observacion}

